%!TEX root = rapport.tex

\chapter{Representation of lambda terms in Python}

A lambda term is represented as a directed acyclic graph (DAG). There are
three types of nodes: abstraction, application and variable.
\begin{description}
	\item[Abstraction nodes] define scopes for bound variables and can be thought
	of as the encoding of the $\lambda$s in the DAG-structure.
	
	\item[Application nodes] have two children. The first of these
	is either the DAG-representation of some \emph{active} subterm of 
	$M$, or another application node. The second child is the DAG-representation
	of the ``argument'', that is, the lambda term that are to be substituted
	in place for a bound variable. An application node represents a possible
	one-step $\beta$-reduction.
	
	\item[Variable nodes] represent free or bound variables. A bound 
	variable is associated with a scope, represented by an ancestral
	abstraction node.
\end{description}
A term that is not in $\beta$-normal-form has at least one $\beta$-contractum,
and therefore the DAG-representation of that term contains at least one
application node. More specifically, for each possible one-step
$\beta$-reduction in a term $M$, its DAG will contain one application node.

The representation is straightforwardly encoded in Python with a general node
class, \pycode{Node}, which has a subclass for each node type. Each
\pycode{Node}-object maintains a list of its immediate descendants. The
abstraction nodes furthermore maintain a list of variable nodes representing
the particular variable that the abstraction node ``binds''. These variable
nodes need not be direct descendants.

% Callback functions? Hvad skal de bruges til?


% \x.xy bliver til
%    L
%    |
%    @
%   / \ 
%  x   y


% application node: @
% abstraction node: L
% variable node: x,y..

% identify each of the possible reducts: find all active subterms.


% (\x.xx)(\y.\z.zyz) == (\x.xx)(\y.\z.(zy)z)
%        @
%       / \
%      L   ------L
%      |         |
%      @         @
%     / \       / \
%    x  x      @   z
%             / \
%            z   y
% .... or:
%       @
%      / \
%    Lx   Ly
%     |   |
%    xx   Lz
%         |
%         zyz

% (\x.xx)(\y.yy)(\z.zz) == ((\x.xx)(\y.yy)) (\z.zz)
%          @
%         / \
%        @   Lz
%       / \  |
%     Lx  Ly zz
%      |   | 
%     xx  yy 
% .... becomes:
%       @
%      / \
%     @   --- Lz
%    / \      |
%   Ly  Ly    zz
%   |   |
%   yy  yy
%

% I(Iz) -> Iz
% (\x.x)((\y.y)z) -> (\y.y)z
%          @
%         / \
%        Lx  ---@
%        |     / \
%        x    Ly  z
%             |
%             y
% indeholder to one-step betareductions:
%     @
%    / \
%   Ly  z
%   |
%   y
% og (ækvivalent!):
%     @
%    / \
%   Lx  z
%   |
%   x




%% Dicts

% For every active subterm n in M, S(M) 